\section{Four-Valued Propositional Kernel}
\label{sec:kernel}

\statusestablished

The propositional base follows the information structure of Belnap--Dunn First-Degree Entailment \cite{Belnap1977,Dunn1976}. Every formula \(\varphi\) receives two independent bits
\[
  v(\varphi)=(t(\varphi),f(\varphi))\in\{0,1\}^2,
\]
where \(t(\varphi)=1\) represents positive support and \(f(\varphi)=1\) negative support. I write \(\pos{\varphi}\) and \(\neginfo{\varphi}\) for these two satisfaction conditions.

\begin{center}
\begin{tabular}{ccc}
\toprule
Value & Pair & Reading \\
\midrule
\(\True\) & \((1,0)\) & positive support only \\
\(\False\) & \((0,1)\) & negative support only \\
\(\Both\) & \((1,1)\) & both positive and negative support \\
\(\Neither\) & \((0,0)\) & neither positive nor negative support \\
\bottomrule
\end{tabular}
\end{center}

Atoms receive positive and negative support independently. The connectives are fixed by the bilateral clauses
\begin{align*}
\pos{\neg\varphi} &\iff \neginfo{\varphi}, &
\neginfo{\neg\varphi} &\iff \pos{\varphi},\\
\pos{(\varphi\land\psi)} &\iff \pos{\varphi}\land\pos{\psi}, &
\neginfo{(\varphi\land\psi)} &\iff \neginfo{\varphi}\lor\neginfo{\psi},\\
\pos{(\varphi\lor\psi)} &\iff \pos{\varphi}\lor\pos{\psi}, &
\neginfo{(\varphi\lor\psi)} &\iff \neginfo{\varphi}\land\neginfo{\psi}.
\end{align*}

Positive semantic consequence preserves positive satisfaction:
\[
  \Gamma\models_{\FDE}\varphi
\]
iff every valuation positively satisfying every member of \(\Gamma\) also positively satisfies \(\varphi\). Equivalently, the positively designated values are \(\{\True,\Both\}\). Negative satisfaction remains separately observable and will be used for bilateral meta-conditions in later sections.

Crucially, no object-language implication is fixed at this stage. Although \(\neg\varphi\lor\psi\) is definable, treating it as \emph{the} conditional would prematurely determine the lifting of A2, D1, D2, and D3.

\begin{proposition}[Paraconsistency]
The kernel is non-explosive: \(p,\neg p\not\models_{\FDE}q\).
\end{proposition}
\begin{proof}
Let \(v(p)=\Both\) and \(v(q)=\Neither\). Then \(p\) and \(\neg p\) are both positively satisfied, while \(q\) is not.
\end{proof}

\begin{proposition}[Paracompleteness]
Excluded middle is not valid: \(\not\models_{\FDE}p\lor\neg p\).
\end{proposition}
\begin{proof}
Let \(v(p)=\Neither\). Negation preserves \(\Neither\), so neither disjunct is positively satisfied and hence neither is their disjunction.
\end{proof}

\begin{proposition}[Classical recovery of the propositional fragment]
If atomic valuations are restricted to \(\{\True,\False\}\), then the clauses for \(\neg\), \(\land\), and \(\lor\) remain inside \(\{\True,\False\}\) and coincide with the classical truth tables.
\end{proposition}
\begin{proof}
Immediate by structural induction on formulas using the bilateral clauses above.
\end{proof}

The full gate specification is maintained in \texttt{docs/FOUR\_VALUED\_KERNEL.md}.
